% Supported v0.5 diversity-statement example.
\documentclass[type=diversity]{careerdossier-statement}
% examples/profiles/profile-academic.tex
%
% Shared profile metadata used by the academic-CV example and the corresponding
% industry-resume smoke fixture. Keeping it independent of document layout
% proves that both document families consume the same \CDossierSetup interface.

\CDossierSetup{
  name     = {Ada Lovelace},
  headline = {Researcher in Analytical Computing},
  email    = {ada.lovelace@example.com},
  location = {Toronto, Ontario, Canada},
  website  = {example.com/ada-lovelace},
  %% The bare Google Scholar identifier, expanded to
  %% scholar.google.com/citations?user=ada-example for both the displayed text
  %% and the link. Pasting the whole address works too and is left as written.
  %% See docs/API.md, "Accepted forms for the web-profile keys".
  scholar  = {ada-example},
  orcid    = {0000-0002-1825-0097},
  affiliation = {Example University},
}

\CDossierStatementSetup{
  application-context={Application for Assistant Professor, Example University},
  application-id={EDIA-2026-17}
}
\begin{document}
\MakeCDossierStatementHeader

\section*{Approach to Inclusive Academic Practice}

I understand equity, diversity, inclusion, and accessibility as responsibilities
expressed through ordinary academic decisions: whose prior knowledge is
recognized, how expectations are communicated, which barriers are treated as
changeable, and how people can influence the environments in which they learn
and work. Inclusive practice is therefore not a separate activity added after a
course or project is designed. It is a method of examining design choices,
seeking evidence about their effects, and revising them with the people most
affected.

My approach joins access with rigor. Students and researchers should encounter
challenging questions, but difficulty should come from the intellectual work,
not from hidden conventions or avoidable obstacles. Clear criteria, accessible
materials, multiple ways to prepare and participate, and timely feedback make
standards more visible. They also make it possible to identify where a system,
rather than an individual's effort, is limiting participation.

\section*{Evidence from Teaching and Mentoring}

In teaching, I use transparent assignment briefs that explain purpose,
prerequisite knowledge, process, and evaluation criteria. Before a major
project, students examine anonymized examples and discuss how the rubric applies
to them. This practice reduces the advantage held by students who already know
how to infer academic expectations and gives the whole class a language for
asking precise questions about quality.

I design participation so that fast public speech is not the only visible form
of engagement. Students may formulate an idea in writing, test it with a partner,
and then bring it to a larger discussion. Collaborative tasks use rotating roles
and shared records, which makes contributions inspectable and prevents one
person from becoming the permanent technical authority. Anonymous check-ins
help me notice patterns that students may not raise publicly, including pace,
accessibility, and group dynamics.

In mentorship, I make the unwritten work of research explicit. New researchers
receive examples of project plans, documentation, review practices, and
authorship conversations. We establish meeting expectations together and revisit
them as the project changes. Rather than assuming that every student has the
same professional network, I discuss how to ask for feedback, prepare for a
conference, evaluate an opportunity, and recover from a rejected submission.

These practices have taught me that a single intervention rarely produces a
complete outcome. Access requires follow-through. When students identify a
barrier, I record the change made, evaluate whether it helped, and preserve the
lesson for future course or project design. This turns inclusion from a promise
into an accountable cycle of action and evidence.

\section*{Research, Service, and Community}

My research groups use written contribution guidelines, review checklists, and
decision records so that collaboration does not depend on proximity to the most
senior person. Documentation gives people entering at different stages a way to
understand why choices were made and to challenge them. I also favour open,
well-described research outputs when ethical and legal conditions permit,
because usable documentation can broaden who is able to examine and extend a
result.

In service, I focus on processes that shape participation: how opportunities
are advertised, how meetings are facilitated, and how work is recognized. I
support advance agendas, accessible remote participation, explicit selection
criteria, and rotation of visible and less visible responsibilities. These are
small structural choices, but together they affect who can contribute and whose
work becomes legible to others.

\section*{Future Contributions at Example University}

At Example University, I would begin by learning from students, staff, faculty,
and existing accessibility and EDIA initiatives rather than assuming that a
portable solution fits every context. I would contribute transparent assessment
resources, structured peer-mentoring practices, and accessible open materials
to courses in computational science. In research supervision, I would establish
shared expectations for communication, credit, documentation, and conflict
resolution at the start of each project.

I would also collaborate with colleagues to examine participation and outcomes
across gateway courses, using appropriate evidence to identify where curriculum
or policy creates avoidable barriers. My commitment is to pair specific action
with reflection and accountability: to state what I can change, invite critique,
measure what happened, and continue learning with the community.

\end{document}
